% Chapter — Rotation kernel (FR-3, D-7). Follows the mandatory FR-29
% six-section template: purpose; assumptions and domain of validity;
% derivation; discretization and implementation; validation evidence;
% references. The equation labels eq:rotations:hamprod,
% eq:rotations:quat2dcm, eq:rotations:r1r2r3, eq:rotations:shepperd,
% eq:rotations:e321, and eq:rotations:e313 are echoed verbatim in comments
% in cpp/src/rotation.cpp and cpp/include/star/rotation.hpp (FR-29
% equation-label-to-code traceability); renaming them breaks that trace.
\chapter{Rotation Kernel}
\label{ch:rotations}

\section{Purpose}
\label{sec:rotations:purpose}

This chapter documents the rotation kernel (FR-3, decision D-7): Hamilton
quaternion algebra, direction-cosine matrices (DCMs), the elementary frame
rotations $\mat{R}_1$, $\mat{R}_2$, $\mat{R}_3$, and the 3-2-1 and 3-1-3
Euler-angle sequences, together with all conversions among these
representations. The kernel is the attitude substrate of the project: the
Earth, Moon, and Mars frame constructions of Chapter~\ref{ch:frames} are
compositions of its elementary rotations, and the Phase~4 attitude dynamics
will integrate quaternions that this kernel converts and normalizes. It is
implemented in \texttt{cpp/src/rotation.cpp} with its interface in
\texttt{cpp/include/star/rotation.hpp}; every convention below restates and
implements the notation chapter's ratification of decision D-7
(Section~\ref{sec:notation:quaternion}).

\section{Assumptions and domain of validity}
\label{sec:rotations:domain}

Assumptions:
\begin{enumerate}
  \item \textbf{Unit quaternions for attitude.} Attitude operations
    (\texttt{dcm\_from\_quat}, \texttt{quat\_transform}) assume unit-norm
    input; \texttt{quat\_normalize} provides the projection and is exact
    to a few ulp. Normalizing a zero or non-finite quaternion would
    fabricate an attitude, so it throws \texttt{std::domain\_error}
    (surfaced in Python as \texttt{ValueError}) instead of defaulting.
  \item \textbf{Proper rotation matrices for extraction.} The DCM-input
    functions (\texttt{quat\_from\_dcm}, the Euler extractions) assume an
    orthonormal matrix with determinant $+1$. Orthonormality is a caller
    invariant, maintained by construction everywhere in the core and gated
    by the property tests (\testid{rotation_dcm_orthonormality_properties});
    it is not re-checked per call inside the deterministic loop.
  \item \textbf{Double-precision conditioning near gimbal lock.} Within an
    angular distance $\delta$ of an Euler sequence's singular middle angle,
    the \emph{individual} first and third angles are ill-conditioned: an
    absolute matrix-element perturbation of order the machine epsilon
    $\varepsilon$ produces angle errors of order $\varepsilon/\delta$,
    because the four elements that determine those angles all carry a
    factor $\cos a_2$ (3-2-1) or $\sin a_2$ (3-1-3) of magnitude $\delta$.
    The recomposed \emph{rotation} degrades with the same $O(\varepsilon/
    \delta)$ envelope, since the two ill-conditioned angle errors arise
    from independent matrix elements and do not cancel. The implementation
    is tested against the conservative envelope $100\,\varepsilon/\delta$
    and observed at $2$--$5\,\varepsilon/\delta$
    (Section~\ref{sec:rotations:validation}).
\end{enumerate}

Domain of validity: all of $SO(3)$ for quaternion and DCM operations; all
of $SO(3)$ for Euler \emph{composition}; Euler \emph{extraction} recovers
individually meaningful angles only away from the sequence's gimbal lock
($|a_2| = \pi/2$ for 3-2-1; $a_2 \in \{0, \pi\}$ for 3-1-3), with the
quantified degradation above inside that neighborhood.

Out-of-domain behavior, matching the code exactly:
\begin{itemize}
  \item A non-orthonormal matrix passed to \texttt{quat\_from\_dcm} or an
    Euler extraction returns the parameters of a nearby rotation without
    detection --- there is no per-call orthonormality check (see
    assumption 2). The guard is upstream: every matrix produced by the
    core is orthonormal to machine precision by construction.
  \item At \emph{exact} gimbal lock (the singular elements identically
    zero) only the sum or difference of the first and third angles is
    observable. The deterministic policy is $a_1 = 0$ with the whole
    degenerate rotation assigned to $a_3$; the recomposed rotation is
    exact. There is no warning --- the policy is a documented convention,
    not an error.
  \item \texttt{quat\_normalize} throws on zero or non-finite input
    (assumption 1); no other kernel function throws.
\end{itemize}

\section{Derivation}
\label{sec:rotations:derivation}

The conventions and closed forms follow Markley and
Crassidis~\cite{markley2014}, transcribed to this project's scalar-first
component ordering; the notation chapter
(Section~\ref{sec:notation:quaternion}) is the authoritative statement of
the conventions themselves.

\subsection{Hamilton product and the attitude matrix}

With scalar-first components $\quat{q} = [\,q_w, \vv{q}_v\,]$,
$\vv{q}_v = [q_x, q_y, q_z]^{\mathsf T}$, the Hamilton product
(equation~\eqref{eq:notation:hamiltonproduct}) written out per component is
\begin{equation}
  \quat{p} \quatmul \quat{q} =
  \begin{bmatrix}
    p_w q_w - p_x q_x - p_y q_y - p_z q_z \\
    p_w q_x + p_x q_w + p_y q_z - p_z q_y \\
    p_w q_y - p_x q_z + p_y q_w + p_z q_x \\
    p_w q_z + p_x q_y - p_y q_x + p_z q_w
  \end{bmatrix}.
  \label{eq:rotations:hamprod}
\end{equation}
For a unit quaternion the inverse is the conjugate
$\quat{q}^{*} = [\,q_w, -\vv{q}_v\,]$. Embedding a vector as the pure
quaternion $\tilde{\vv{v}} = [\,0, \vv{v}\,]$ and expanding the frame
transformation sandwich $\tilde{\vv{v}}^{B} = \quat{q}^{-1} \quatmul
\tilde{\vv{v}}^{I} \quatmul \quat{q}$ of
equation~\eqref{eq:notation:sandwich} by two applications of
equation~\eqref{eq:rotations:hamprod} gives, after collecting terms in
$q_w^2$, $q_w \vv{q}_v$, and $\vv{q}_v \vv{q}_v^{\mathsf T}$,
\begin{equation}
  \dcm{A}{B}(\quat{q})
  = \bigl(q_w^{2} - \vv{q}_v \cdot \vv{q}_v\bigr) \mat{I}_3
    + 2\, \vv{q}_v \vv{q}_v^{\mathsf T}
    - 2\, q_w \skewm{\vv{q}_v}
  =
  \begin{bmatrix}
    q_w^2{+}q_x^2{-}q_y^2{-}q_z^2 & 2(q_x q_y + q_w q_z) & 2(q_x q_z - q_w q_y) \\
    2(q_x q_y - q_w q_z) & q_w^2{-}q_x^2{+}q_y^2{-}q_z^2 & 2(q_y q_z + q_w q_x) \\
    2(q_x q_z + q_w q_y) & 2(q_y q_z - q_w q_x) & q_w^2{-}q_x^2{-}q_y^2{+}q_z^2
  \end{bmatrix},
  \label{eq:rotations:quat2dcm}
\end{equation}
the element form of the attitude matrix of
equation~\eqref{eq:notation:quat2dcm}~\cite{markley2014}. Because
$\quat{q}$ and $-\quat{q}$ produce the same matrix, every extraction below
resolves the sign to $q_w \geq 0$.

\subsection{Elementary frame rotations}

The elementary rotations are \emph{frame} rotations (coordinate
transformations): $\mat{R}_3(\theta)$ maps the coordinates of a fixed
vector into a frame rotated by $+\theta$ about the $+z$ axis,
\begin{equation}
  \mat{R}_1(\theta) =
  \begin{bmatrix}
    1 & 0 & 0 \\ 0 & \cos\theta & \sin\theta \\ 0 & -\sin\theta & \cos\theta
  \end{bmatrix},
  \;
  \mat{R}_2(\theta) =
  \begin{bmatrix}
    \cos\theta & 0 & -\sin\theta \\ 0 & 1 & 0 \\ \sin\theta & 0 & \cos\theta
  \end{bmatrix},
  \;
  \mat{R}_3(\theta) =
  \begin{bmatrix}
    \cos\theta & \sin\theta & 0 \\ -\sin\theta & \cos\theta & 0 \\ 0 & 0 & 1
  \end{bmatrix}.
  \label{eq:rotations:r1r2r3}
\end{equation}
Each is the DCM of the axis-angle quaternion $[\cos(\theta/2),
\sin(\theta/2)\,\vv{e}_i]$ under equation~\eqref{eq:rotations:quat2dcm}
(verified as a hand-value test), and the signs coincide with the
$\mathtt{Rx}/\mathtt{Ry}/\mathtt{Rz}$ primitives of the IAU SOFA
vector-matrix library~\cite{erfa201}, which the Earth-orientation chain of
Chapter~\ref{ch:frames} composes.

\subsection{Euler sequences}

Angles are stated in \emph{application order}: the sequence name lists the
rotation axes in the order applied. The two sequences the project requires
(aerospace 3-2-1; the 3-1-3 of the lunar libration and Mars pole
constructions) are
\begin{align}
  \dcm{A}{B} &= \mat{R}_1(a_3)\,\mat{R}_2(a_2)\,\mat{R}_3(a_1)
  \nonumber \\
  &=
  \begin{bmatrix}
    c_2 c_1 & c_2 s_1 & -s_2 \\
    s_3 s_2 c_1 - c_3 s_1 & c_3 c_1 + s_3 s_2 s_1 & s_3 c_2 \\
    s_3 s_1 + c_3 s_2 c_1 & c_3 s_2 s_1 - s_3 c_1 & c_3 c_2
  \end{bmatrix}
  \qquad \text{(3-2-1)},
  \label{eq:rotations:e321} \\[2pt]
  \dcm{A}{B} &= \mat{R}_3(a_3)\,\mat{R}_1(a_2)\,\mat{R}_3(a_1)
  \nonumber \\
  &=
  \begin{bmatrix}
    c_3 c_1 - s_3 c_2 s_1 & c_3 s_1 + s_3 c_2 c_1 & s_3 s_2 \\
    -s_3 c_1 - c_3 c_2 s_1 & c_3 c_2 c_1 - s_3 s_1 & c_3 s_2 \\
    s_2 s_1 & -s_2 c_1 & c_2
  \end{bmatrix}
  \qquad \text{(3-1-3)},
  \label{eq:rotations:e313}
\end{align}
with $c_i = \cos a_i$, $s_i = \sin a_i$, obtained by direct multiplication
of equation~\eqref{eq:rotations:r1r2r3}; the sequence definitions follow
Markley and Crassidis~\cite{markley2014}.

Extraction inverts the element identities. For 3-2-1, from
equation~\eqref{eq:rotations:e321}: $C_{13} = -\sin a_2$ and
$C_{11}^2 + C_{12}^2 = \cos^2 a_2$, so
\begin{equation}
  a_2 = \operatorname{atan2}\!\Bigl(-C_{13},\, \sqrt{C_{11}^2 + C_{12}^2}
  \Bigr),
  \qquad
  a_1 = \operatorname{atan2}(C_{12}, C_{11}),
  \qquad
  a_3 = \operatorname{atan2}(C_{23}, C_{33}),
  \label{eq:rotations:e321extract}
\end{equation}
which places $a_2 \in [-\pi/2, \pi/2]$. The two-argument arctangent of the
$(\sin, \cos)$ pair is used for $a_2$ rather than $\arcsin(-C_{13})$
because the derivative of $\arcsin$ diverges at $\pm 1$: near lock the
$\arcsin$ form loses half the significand, while the pair form keeps $a_2$
accurate to machine precision arbitrarily close to lock. For 3-1-3, from
equation~\eqref{eq:rotations:e313}: $C_{33} = \cos a_2$,
$C_{13}^2 + C_{23}^2 = \sin^2 a_2$, and
\begin{equation}
  a_2 = \operatorname{atan2}\!\Bigl(\sqrt{C_{13}^2 + C_{23}^2},\, C_{33}
  \Bigr),
  \qquad
  a_1 = \operatorname{atan2}(C_{31}, -C_{32}),
  \qquad
  a_3 = \operatorname{atan2}(C_{13}, C_{23}),
  \label{eq:rotations:e313extract}
\end{equation}
with $a_2 \in [0, \pi]$. At exact lock the arguments of the $a_1$ and
$a_3$ arctangents are all identically zero and the policy of
Section~\ref{sec:rotations:domain} applies: $a_1 = 0$, and $a_3$ is read
from the sub-block that then depends only on the observable combination
(for 3-2-1 at $a_2 = +\pi/2$, $C_{21} = \sin(a_3 - a_1)$ and
$C_{22} = \cos(a_3 - a_1)$; the other lock branches are analogous and
spelled out in the source).

\subsection{Robust quaternion extraction (Shepperd's method)}

Taking traces and diagonal elements of
equation~\eqref{eq:rotations:quat2dcm} with $\norm{\quat{q}} = 1$ gives
\begin{equation}
  1 + \operatorname{tr}\dcm{A}{B} = 4 q_w^2,
  \qquad
  1 + 2 C_{ii} - \operatorname{tr}\dcm{A}{B} = 4 q_i^2
  \quad (i = x, y, z),
  \label{eq:rotations:shepperd}
\end{equation}
while the off-diagonal sums and differences yield the mixed products, e.g.
$C_{23} - C_{32} = 4 q_w q_x$ and $C_{12} + C_{21} = 4 q_x q_y$. Since the
four quantities in equation~\eqref{eq:rotations:shepperd} sum to $4$, the
largest is at least $1$; Shepperd's method~\cite{shepperd1978} computes
that component by square root and the remaining three by division with the
guaranteed-large denominator, so the extraction never divides by a small
quantity and is numerically robust for every attitude --- unlike the
trace-only formula, which cancels catastrophically near $180^\circ$
rotations. The four golden cases \texttt{shepperd\_*\_dominant} make each
component dominant in turn, exercising all four branches.

\subsection{Conditioning near gimbal lock}

Let the input matrix carry absolute element perturbations of order
$\varepsilon$ (the rounding floor of any upstream computation, e.g.
equation~\eqref{eq:rotations:quat2dcm}). At angular distance $\delta$ from
lock, the elements determining $a_1$ and $a_3$ have magnitude
$O(\delta)$ --- each carries the factor $\cos a_2$ or $\sin a_2$ --- so the
extracted angles err by $O(\varepsilon/\delta)$. Recomposition does not
cancel these errors: the 3-2-1 element $C_{22} = c_3 c_1 + s_3 s_2 s_1$
depends on the error combination $(e_1 - e_3)\sin(a_3 - a_1)$ at lock, and
$e_1$, $e_3$ derive from disjoint element sets. The end-to-end round-trip
error is therefore bounded by $K \varepsilon / \delta$ with $K$ of order
unity; the validation evidence tests $K = 100$ as a conservative envelope
and observes $K \approx 2$--$5$
(Table~\ref{tab:rotations:validation-core}). For uniformly random
attitudes the probability of falling within $10^{-3}$ of lock is below
$10^{-6}$ per draw, and the Phase~2 criterion-6a bound of
\qty{1e-13}{\radian} holds with two orders of margin over the observed
worst case.

\section{Discretization and implementation}
\label{sec:rotations:implementation}

\begin{enumerate}
  \item \textbf{Module and traceability.} The kernel lives in
    \texttt{cpp/src/rotation.cpp}; the source comments echo the labels
    \texttt{eq:rotations:hamprod}, \texttt{eq:rotations:quat2dcm},
    \texttt{eq:rotations:r1r2r3}, \texttt{eq:rotations:shepperd},
    \texttt{eq:rotations:e321}, and \texttt{eq:rotations:e313} verbatim.
  \item \textbf{Eigen storage mapping (D-7).} Quaternions are stored as
    \texttt{Eigen::Quaterniond}, whose constructor is scalar-first while
    its internal \texttt{coeffs()} storage is scalar-last; the mandatory
    mapping table and its two binding rules are in the notation chapter
    (Section~\ref{sec:notation:eigen}). The binding layer emits
    quaternions scalar-first through the named accessors exclusively.
    Eigen's \texttt{.toRotationMatrix()} implements the \emph{active}
    convention and would return the transpose of
    equation~\eqref{eq:rotations:quat2dcm}; core code uses this kernel's
    \texttt{dcm\_from\_quat} for frame transformations, never
    \texttt{.toRotationMatrix()}.
  \item \textbf{One code path for vector transformation.}
    \texttt{quat\_transform} evaluates the sandwich through
    \texttt{dcm\_from\_quat}, so the DCM and the quaternion action on
    vectors agree structurally rather than by test alone.
  \item \textbf{Explicit Hamilton product.} \texttt{quat\_multiply} writes
    out equation~\eqref{eq:rotations:hamprod} component by component
    rather than delegating to Eigen's \texttt{operator*}, so the
    convention is visible at the code and immune to any upstream semantic
    change; the golden tests pin the equivalence.
  \item \textbf{Sign conventions.} \texttt{quat\_from\_dcm} returns
    $q_w \geq 0$; extractions are therefore reproducible bit for bit
    (D-10). All compositions use fixed evaluation order; there is no
    fast-math and no contraction (D-10 build flags).
  \item \textbf{Exact-lock branch.} The Euler extractions branch to the
    lock policy only when the singular elements are \emph{exactly} zero
    (the hypot of two doubles vanishes only then); every other input takes
    the general branch, whose conditioning is quantified above.
\end{enumerate}

\section{Validation evidence}
\label{sec:rotations:validation}

Tables~\ref{tab:rotations:validation-core}
and~\ref{tab:rotations:validation-py} list the tests that constitute this
chapter's validation evidence. Each test identifier is machine-checked to
exist in the test tree by \texttt{scripts/lint\_chapters.py}; golden
provenance and tolerances are recorded in
\texttt{tests/golden/rotations/manifest.toml}.

\begin{table}[htbp]
  \centering
  \caption{Validation evidence for the rotation kernel: core (doctest,
    \texttt{cpp/tests/}).}
  \label{tab:rotations:validation-core}
  \begin{tabular}{lp{8.6cm}}
    \toprule
    Test ID & Acceptance basis \\
    \midrule
    \testid{rotation_quat_dcm_golden} &
      Equation~\eqref{eq:rotations:quat2dcm} vs.\ doubly constructed
      golden matrices (ERFA primitives and independent NumPy closed form)
      to $\leq 10^{-15}$ per element; Shepperd extraction recovers the
      committed quaternion to $\leq 10^{-15}$ per component with
      $q_w \geq 0$, all four pivot branches exercised. \\
    \testid{rotation_euler_golden} &
      Equations~\eqref{eq:rotations:e321}
      and~\eqref{eq:rotations:e313} vs.\ ERFA-composed goldens to
      $\leq 10^{-15}$ per element, including gimbal-lock-adjacent angle
      sets. \\
    \testid{rotation_primitive_axes_hand_values} &
      Equation~\eqref{eq:rotations:r1r2r3} against hand values at
      $\pm\pi/2$ and $\pi/6$ and against the equivalent axis-angle
      quaternions. \\
    \testid{rotation_quat_algebra_properties} &
      Unit norm after normalization; $(\quat{p} \quatmul \quat{q})^{*} =
      \quat{q}^{*} \quatmul \quat{p}^{*}$; $\quat{p} \quatmul
      \quat{p}^{*} = $ identity; the composition rule of
      equation~\eqref{eq:notation:quatcomp}; sandwich--DCM agreement. \\
    \testid{rotation_dcm_orthonormality_properties} &
      $\mat{C}^{\mathsf T}\mat{C} = \mat{I}_3$ and $\det \mat{C} = +1$ to
      $2 \times 10^{-15}$ (a few ulp) for every produced DCM: 200 random
      attitudes, Euler compositions, primitive products. \\
    \testid{rotation_quat_dcm_euler_roundtrip} &
      Phase~2 exit criterion 6a: quaternion$\to$DCM$\to$Euler round trips
      recover the input to $\leq \qty{1e-13}{\radian}$ rotation-equivalent
      over 1{,}000 seeded random attitudes (observed worst
      \qty{1.3e-15}{\radian}) and near-lock draws at $\delta = 10^{-1},
      10^{-2}$; degradation at $\delta = 10^{-4}, 10^{-6}, 10^{-8}$ held
      to the $100\,\varepsilon/\delta$ envelope (observed \num{5.4e-13},
      \num{1.1e-10}, and \qty{1.7e-9}{\radian}); exact-lock
      policy recomposes to $\leq 5 \times 10^{-16}$ per element. \\
    \bottomrule
  \end{tabular}
\end{table}

\begin{table}[htbp]
  \centering
  \caption{Validation evidence for the rotation kernel: binding surface
    (pytest, \texttt{tests/python/}).}
  \label{tab:rotations:validation-py}
  \begin{tabular}{lp{7.2cm}}
    \toprule
    Test ID & Acceptance basis \\
    \midrule
    \testid{test_rotation_bindings_match_goldens} &
      The pybind11 surface reproduces the quaternion and Euler goldens
      (scalar-first tuples; $\leq 10^{-15}$ per element). \\
    \testid{test_rotation_bindings_roundtrip} &
      Criterion-6a round trips through the bindings on the identical
      seeded RNG stream the doctest consumes; \texttt{ValueError} for a
      degenerate quaternion. \\
    \bottomrule
  \end{tabular}
\end{table}

\section{References}
\label{sec:rotations:references}

This chapter relies on Markley and Crassidis~\cite{markley2014} for the
quaternion conventions, the attitude-matrix closed form, and the Euler
sequence definitions; on Shepperd~\cite{shepperd1978} for the robust
DCM-to-quaternion extraction; on the ERFA/SOFA vector-matrix
primitives~\cite{erfa201} for the frame-rotation sign convention the
golden vectors are built with; and on Eigen~\cite{eigenweb} for the
storage-order facts restated from the notation chapter. Full bibliographic
details appear in the bibliography.
